\clearpage\pagestyle{otMain}\setlength{\otBlocEspacement}{2.5mm}\fontsize{11}{15}\selectfont\corrigesChapitre{11}{Longueurs et périmètres}
\begin{blocLivre}{cor-11-1}\begin{corrige}{1}{La bonne unité}\demarche{Associer chaque objet à une référence connue.}
Épaisseur : $\res{\text{mm}}$. Cahier : $\res{\text{cm}}$. Immeuble : $\res{\text m}$. Distance entre villes : $\res{\text{km}}$.
\end{corrige}\end{blocLivre}
\begin{blocLivre}{cor-11-2}\begin{corrige}{2}{Convertir en raisonnant}\demarche{Citer la relation entre les deux unités.}
$1$ m $=100$ cm, donc $3{,}5$ m $=\res{350}$ cm. $1$ km $=1\,000$ m, donc 2 km $=\res{2\,000}$ m. $10$ mm $=1$ cm, donc 48 mm $=\res{4{,}8}$ cm.
\end{corrige}\end{blocLivre}
\begin{blocLivre}{cor-11-3}\begin{corrige}{3}{Comparer des longueurs}\demarche{Convertir dans la même unité avant de comparer.}
$0{,}5$ m $=50$ cm et $50>48$ : $\res{0{,}5\ \mathrm m>48\ \mathrm{cm}}$. $1{,}5$ km $=1\,500$ m : $\res{1\,200\ \mathrm m<1{,}5\ \mathrm{km}}$.
\end{corrige}\end{blocLivre}
\begin{blocLivre}{cor-11-4}\begin{corrige}{4}{Estimer sans mesurer}\demarche{Comparer mentalement à des objets connus.}
La gomme mesure environ $\res{4\ \mathrm{cm}}$, le terrain $\res{100\ \mathrm m}$, la chaise $\res{90\ \mathrm{cm}}$. Les références évitent les réponses invraisemblables.
\end{corrige}\end{blocLivre}
\begin{blocLivre}{cor-11-5}\begin{corrige}{5}{Le tour du quadrilatère}\demarche{Additionner tous les côtés du contour.}
On additionne les quatre côtés : $3+4{,}5+2{,}5+4=\res{14}$ cm. Regrouper $4{,}5+2{,}5=7$ rend le calcul rapide.
\label{fig-cor-11-5}\figureGeo[corrige]{per-5}
\end{corrige}\end{blocLivre}
\begin{blocLivre}{cor-11-6}\begin{corrige}{6}{Le carré}\demarche{Multiplier le côté par 4.}
Les quatre côtés sont égaux : $4\times7=\res{28}$ cm.
\label{fig-cor-11-6}\figureGeo[corrige]{per-6}
\end{corrige}\end{blocLivre}
\begin{blocLivre}{cor-11-7}\begin{corrige}{7}{Le rectangle}\demarche{Doubler la somme longueur plus largeur.}
Longueur plus largeur : $8+3{,}5=11{,}5$ cm, puis le double : $\res{23}$ cm.
\label{fig-cor-11-7}\figureGeo[corrige]{per-7}
\end{corrige}\end{blocLivre}
\begin{blocLivre}{cor-11-8}\begin{corrige}{8}{Retrouver le côté}\demarche{Partager le périmètre en quatre.}
Le côté est le quart du périmètre : $36:4=\res9$ cm. Contrôle : $4\times9=36$.
\end{corrige}\end{blocLivre}
\begin{blocLivre}{cor-11-9}\begin{corrige}{9}{Des kilomètres aux mètres}\demarche{Multiplier ou partager par la relation entre unités.}
$1{,}25\times1\,000=\res{1\,250}$ m. Cent centimètres font un mètre : 230 cm $=\res{2{,}30}$ m, c’est-à-dire $2{,}3$ m.
\end{corrige}\end{blocLivre}
\begin{blocLivre}{cor-11-10}\begin{corrige}{10}{La figure en L}\demarche{Déduire les côtés manquants avant d’additionner.}
Les deux côtés sans mesure se déduisent : le petit côté du haut vaut $8-5=3$ cm et le côté vertical intérieur vaut $6-2=4$ cm. Le contour vaut $8+6+5+4+3+2=\res{28}$ cm.
\label{fig-cor-11-10}\figureGeo[corrige]{per-10}
\end{corrige}\end{blocLivre}
\begin{blocLivre}{cor-11-11}\begin{corrige}{11}{La clôture du jardin}\demarche{Retirer l’ouverture après le calcul du tour.}
Périmètre : $2\times(12{,}5+8)=2\times20{,}5=\res{41}$ m. En retirant le portail : $41-1{,}5=\res{39{,}5}$ m de grillage.
\end{corrige}\end{blocLivre}
\begin{blocLivre}{cor-11-12}\begin{corrige}{12}{Les tours de piste}\demarche{Multiplier la longueur d’un tour, puis convertir.}
$400\times5=\res{2\,000}$ m. Comme $1\,000$ m font 1 km : $\res2$ km.
\end{corrige}\end{blocLivre}
\begin{blocLivre}{cor-11-13}\begin{corrige}{13}{Deux chemins sur quadrillage}\demarche{Compter les déplacements horizontaux et verticaux.}
Le rectangle a pour périmètre $2\times(5+3)=16$ cm. En suivant l’escalier, chaque marche remplace un coin sans changer le total horizontal ni vertical : son périmètre vaut aussi $\res{16}$ cm. Deux figures différentes peuvent avoir $\res{\text{le même périmètre}}$.
\label{fig-cor-11-13}\figureGeo[corrige]{per-13}
\end{corrige}\end{blocLivre}
\begin{blocLivre}{cor-11-14}\begin{corrige}{14}{Le cadre de la photo}\demarche{Tout convertir en centimètres avant de conclure.}
Le contour vaut $2\times(15+10)=50$ cm. Le ruban mesure $0{,}45$ m $=45$ cm. Comme $45<50$, il manque $\res5$ cm : le ruban $\res{\text{ne suffit pas}}$.
\end{corrige}\end{blocLivre}
\begin{blocLivre}{cor-11-15}\begin{corrige}{15}{Même périmètre, même figure ?}\demarche{Chercher deux figures différentes de même tour.}
Un carré de côté 4 cm a un périmètre de 16 cm ; un rectangle de 6 cm sur 2 cm aussi : $2\times(6+2)=16$. Les figures sont différentes : l’affirmation est $\res{\text{fausse}}$.
\end{corrige}\end{blocLivre}
